\section*{Capítulo \ref{ch:g12:internal-heat} --- O calor interno da Terra}

\begin{solution}{exo:g12:internal-heat:1}
(a) o \omterm{def:g12:internal-heat:heat-flow}{fluxo de calor}; (b) a \omterm{def:g12:internal-heat:heat-production}{produção de calor radiogênico}; (c) o
\omterm{def:g12:internal-heat:secular-cooling}{resfriamento secular}; (d) a \omterm{def:g12:internal-heat:convection}{convecção térmica}; (e) a \omterm{def:g12:internal-heat:geothermal}{energia geotérmica}.
\end{solution}

\begin{solution}{exo:g12:internal-heat:2}
(1) \qty{0.0916}{W/m^2}. (2) \qty{4.67e13}{W}. (3) \qty{0.029}{K/m}: uma
diferença de um grau Celsius é uma diferença de um kelvin.
\end{solution}

\begin{solution}{exo:g12:internal-heat:3}
Tório-232, urânio-238, potássio-40 e urânio-235. O KamLAND só vê os
geoneutrinos do tório-232 e do urânio-238.
\end{solution}

\begin{solution}{exo:g12:internal-heat:4}
$490/\sqrt{1} = \qty{490}{mW/m^2}$; $490/2 = \qty{245}{mW/m^2}$;
$490/5 = \qty{98}{mW/m^2}$; $490/7 = \qty{70}{mW/m^2}$, cerca da média dos
continentes.
\end{solution}

\begin{solution}{exo:g12:internal-heat:5}
Perto da superfície, a temperatura do solo acompanha os dias e as estações;
só abaixo de alguns metros ela fica na temperatura média do ano, e só mais
fundo ainda ela mostra o calor que vem de baixo. A perfuração agita a rocha
e a água do furo e muda a sua temperatura, por isso se espera até que o
furo tenha voltado à temperatura da rocha.
\end{solution}

\begin{solution}{exo:g12:internal-heat:6}
(1) $(265 - 10)/9.1 \approx \qty{28.0}{\celsius}$ por quilômetro, isto é,
\qty{0.0280}{K/m}. (2) $q = 2.72 \times 0.0280 \approx \qty{0.0762}{W/m^2}$,
cerca de \qty{76}{mW/m^2}. (3) Um pouco acima da média dos continentes,
\qty{70.9}{mW/m^2}, em cerca de \qty{7}{\percent}: um valor continental
típico.
\end{solution}

\begin{solution}{exo:g12:internal-heat:7}
Continentes: $0.0709 \times \num{2.073e14} \approx \qty{1.47e13}{W} =
\qty{14.7}{TW}$. Oceanos: $0.1054 \times \num{3.028e14} \approx
\qty{3.19e13}{W} = \qty{31.9}{TW}$. Total de cerca de \qty{46.6}{TW}, dos
quais $31.9/46.6 \approx \qty{68}{\percent}$, dois terços, saem através do
fundo oceânico.
\end{solution}

\begin{solution}{exo:g12:internal-heat:8}
(1) $\num{1.88e-6} \times 8000 \approx \qty{0.0150}{W/m^2} = \qty{15.0}{mW/m^2}$.
(2) $15.0/70.9 \approx \qty{21}{\percent}$, um quinto do \omterm{def:g12:internal-heat:heat-flow}{fluxo de calor}.
\end{solution}

\begin{solution}{exo:g12:internal-heat:9}
$\sqrt{t} = 490/70.9 \approx 6.91$, logo $t \approx 47.8$ milhões de anos: o
fundo oceânico com mais de cerca de 48 milhões de anos perde menos calor que
o continente médio.
\end{solution}

\begin{solution}{exo:g12:internal-heat:10}
Urânio-235: $2^{1.4/0.704} = 2^{1.99} \approx 4$ vezes mais. Potássio-40:
$2^{2.5/1.248} = 2^{2.00} \approx 4$ vezes mais. Os dois estão duas
meias-vidas atrás.
\end{solution}

\begin{solution}{exo:g12:internal-heat:11}
(1) $10^2/10^{-6} = \qty{1e8}{s}$, cerca de 3 anos. (2) $(10^5)^2/10^{-6} =
\qty{1e16}{s}$, cerca de 320 milhões de anos. (3) $(\num{2.891e6})^2/10^{-6}
\approx \qty{8.4e18}{s}$, cerca de 265 bilhões de anos. O tempo cresce como
o quadrado da distância: dez vezes mais espesso, cem vezes mais longo.
\end{solution}

\begin{solution}{exo:g12:internal-heat:12}
Multiplique cada um por $2^{3/T}$: tório-232, $10.4 \times 2^{3/14.0} \approx
\qty{12.1}{TW}$; urânio-238, $9.6 \times 2^{3/4.468} \approx \qty{15.3}{TW}$;
potássio-40, $4.0 \times 2^{3/1.248} \approx \qty{21.2}{TW}$; urânio-235,
$0.42 \times 2^{3/0.704} \approx \qty{8.1}{TW}$. Total de cerca de
\qty{56.7}{TW}, mais que o dobro dos \qty{24.4}{TW} de hoje. Os dois
isótopos de vida curta dão
$(21.2 + 8.1)/56.7 \approx \qty{52}{\percent}$ dele, contra
$4.42/24.4 \approx \qty{18}{\percent}$ hoje.
\end{solution}

\begin{solution}{exo:g12:internal-heat:13}
A potência é a massa por segundo vezes o calor específico vezes o
aquecimento: $\num{3e8} = 1640 \times 4180 \times \Delta T$, logo $\Delta T =
\num{3e8}/(1640 \times 4180) \approx \qty{44}{K}$. O continente médio perde
\qty{0.0709}{W} por metro quadrado, então \qty{300}{MW} são a perda natural
de calor de $\num{3e8}/0.0709 \approx \qty{4.2e9}{m^2}$, cerca de
\qty{4200}{km^2}: a usina reúne o calor que a Terra deixa sair naturalmente
em milhares de quilômetros quadrados.
\end{solution}

\begin{solution}{exo:g12:internal-heat:14}
(1) 100 milhões de anos é $\num{2.6e11}/\num{1e8} = 2600$ vezes mais curto
que 260 bilhões de anos; 200 milhões de anos, 1300 vezes mais curto.
(2) A convecção leva o calor com a rocha que se move: o tempo para
atravessar o manto é a sua espessura dividida pela velocidade,
$\qty{2891}{km}$ a alguns centímetros por ano, cerca de cem milhões de anos.
A condução não tem esse veículo: o seu tempo cresce como o quadrado da
distância, e ao longo de milhares de quilômetros ele se torna enorme.
\end{solution}

\begin{solution}{exo:g12:internal-heat:15}
(1) $10/22.3 \approx \qty{45}{\percent}$; o resto é o resfriamento do próprio
manto, que perde o seu próprio calor armazenado. (2) O limite fica no raio
$6371 - 2891 =
\qty{3480}{km}$; a sua área é $4\pi (\num{3.48e6})^2 \approx
\qty{1.52e14}{m^2}$, e $\num{1e13}/\num{1.52e14} \approx
\qty{0.066}{W/m^2}$, cerca de \qty{66}{mW/m^2}, uns \qty{72}{\percent} da
média na superfície. (3) O calor atravessa o limite por condução até a
camada mais baixa do manto; aquecida, essa rocha fica mais leve que o manto
acima dela e pode subir, levando o calor para cima: é assim que se pensa
que as plumas começam, em grande profundidade.
\end{solution}

\begin{solution}{pb:g12:internal-heat:1}
\textbf{1.} $4\pi (\num{6.371e6})^2 \approx \qty{5.101e14}{m^2}$, e
$2.073 + 3.028 = 5.101$: sim.

\textbf{2.} $3.028/5.101 \approx \qty{59}{\percent}$, cerca de três quintos.

\textbf{3.} O fundo oceânico é jovem: ele é produzido quente nas dorsais e
esfria à medida que se afasta, e a \omterm{def:g11:lithosphere-asthenosphere:lithosphere}{litosfera} jovem perde muito calor.

\textbf{4.} $490/\sqrt{4} = \qty{245}{mW/m^2}$.

\textbf{5.} $2.72 \times 0.026 \approx \qty{0.0707}{W/m^2} =
\qty{70.7}{mW/m^2}$: um pouco abaixo da média, praticamente igual a ela.

\textbf{6.} $0.0709 \times \num{2.073e14} \approx \qty{14.7}{TW}$.

\textbf{7.} $0.1054 \times \num{3.028e14} \approx \qty{31.9}{TW}$.

\textbf{8.} $14.7 + 31.9 \approx \qty{46.6}{TW}$, os \qty{47}{TW} da
estimativa publicada (a média global do capítulo dá \qty{46.7}{TW}, o valor
usado a seguir).

\textbf{9.} $\num{4.66e13}/\num{5.101e14} \approx \qty{0.0914}{W/m^2}$,
cerca de \qty{91}{mW/m^2}.

\textbf{10.} $0.0243 \times \num{2.073e14} \approx \qty{5.0}{TW}$, um terço
dos \qty{14.7}{TW} perdidos através dos continentes: a radioatividade da
\omterm{def:g12:continental-crust:continental-crust}{crosta continental} é uma grande parte do seu \omterm{def:g12:internal-heat:heat-flow}{fluxo de calor}.

\textbf{11.} Para cada quilograma de urânio, a Terra contém \qty{3.9}{kg}
de tório, que liberam $3.9 \times 26.3 \approx \qty{102.6}{\micro W}$,
contra \qty{94.4}{\micro W} para o urânio-238. O urânio-238 produz assim
$94.4/(94.4 + 102.6) \approx 0.479$ dos \qty{20.0}{TW}: cerca de
\qty{9.6}{TW}. Urânio-235: $9.59 \times 4.10/94.4 \approx \qty{0.42}{TW}$.

\textbf{12.} $20.0 + 0.42 + 4 \approx \qty{24.4}{TW}$.

\textbf{13.} O granito e as outras rochas da \omterm{def:g12:continental-crust:continental-crust}{crosta continental}, cerca de
vinte vezes mais ricos que as rochas ultramáficas do manto.

\textbf{14.} Maior: cada isótopo era mais abundante então, o dobro para cada
meia-vida atrás; há 3 bilhões de anos, o total era mais que o dobro do de
hoje.

\textbf{15.} $46.7 - 24.4 = \qty{22.3}{TW}$.

\textbf{16.} Do calor armazenado desde que a Terra se formou: o
\omterm{def:g12:internal-heat:secular-cooling}{resfriamento secular} do manto e do núcleo.

\textbf{17.} $22.3/46.7 \approx \qty{48}{\percent}$.

\textbf{18.} De $11/46.7 \approx \qty{24}{\percent}$ a $38/46.7 \approx
\qty{81}{\percent}$: os geoneutrinos estreitam muito essa faixa.

\textbf{19.} A \omterm{def:g12:internal-heat:convection}{convecção do manto}, que leva a maior parte do calor; só as
últimas dezenas de quilômetros são atravessadas por condução.

\textbf{20.} $24.4/46.7 \approx 0.52$: a radioatividade fornece
\textbf{cerca da metade} da perda de calor da Terra hoje, e o lento
resfriamento do planeta a outra metade.
\end{solution}
